\chapter{Matrix modes without determinants}\label{ch:matrixmodes}

\section{Characteristic modes are not transfer poles}

Version 0.8.0 adds \code{cmodes}, which searches a square analytic matrix $K(s)$ directly. A characteristic value
$\lambda$ satisfies $K(\lambda)v=0$ for some nonzero $v$. For a structured
model $G=C K^{-1}B+D$, such a value need not be visible as a transfer pole:
the input may not excite it or the output may not observe it. The new
function deliberately does not perform that visibility analysis.

\begin{lstlisting}
K = @(s) [s+1 .2*exp(-s); 0 s+2];
Kp = @(s) [1 -.2*exp(-s); 0 1];
box = [-3 0 -1 1];
[r,info] = cmodes(K,box,'AssumeAnalytic',true, ...
    'Derivative',Kp);
M = cdyn(K,[1;0],[1 0],0,'Dimensions',[2 1 1]);
[m,im] = cmodes(M,box,'AssumeAnalytic',true);
% Both searches return -1 and -2, but G(s) = 1/(s+1).
\end{lstlisting}

\code{AssumeAnalytic=true} asserts that $K$ is regular and analytic in a
neighborhood of the closed rectangle. The boundary must be nonsingular.
A \code{cdyn} argument supplies only its $K$ factor and domain guard;
the other factors are not evaluated. Direct handles must return a
fixed-size square matrix. Sparse inputs are accepted, but the initial
engine uses dense LU/SVD and is intended for small to moderate matrices.

\section{Counting without a product of pivots}

Jacobi's identity connects the matrix and scalar argument principles:
\begin{equation}
 N=\frac{1}{2\pi\ii}\oint_\Gamma
       \operatorname{tr}\bigl(K(s)^{-1}K'(s)\bigr)\,ds
   =\frac{1}{2\pi}\Delta_\Gamma\arg\det K(s).
\end{equation}
Choose diagonal nonsingular scalings $L_0,R_0$ at a regular reference
point and freeze them. Then $A=L_0KR_0$ has the same characteristic values.
At each contour node compute $PA=LU$. The logarithmic determinant
magnitude is $\sum_j\log|U_{jj}|$, and its phase is the sum of the pivot
phases plus the permutation phase. There is no multiplication of pivots.
Neither the changing pivot choices nor singular values are treated as
analytic functions of $s$.

\begin{algobox}{Bounded contour search for matrix modes}
\begin{enumerate}
 \item Select fixed scaling at an admissible regular reference node.
 \item Sample the counterclockwise rectangle. Refine until phase/log-magnitude
       increments are resolved and integer winding agrees twice successively.
 \item If an analytic derivative is supplied, cross-check the outer count
       by integrating $\operatorname{tr}(A\backslash A')$.
 \item Refine candidate locations by a damped bordered Newton solve;
       multiple-count cells can also use the logarithmic derivative.
 \item Confirm a candidate by a small local contour. Compare its algebraic
       count with numerical nullity and the projected derivative.
 \item Otherwise split the cell, shifting internal split lines if needed.
       Accept children only when their counts sum to the parent count.
 \item Report unresolved cells/clusters when any work or resolution limit
       is reached. Never replace failure with a successful empty spectrum.
\end{enumerate}
\end{algobox}

The bordered simple-mode step solves
\begin{equation}
 \begin{bmatrix}A(\lambda)&A'(\lambda)v\\v^*&0\end{bmatrix}
 \begin{bmatrix}\delta v\\\delta\lambda\end{bmatrix}
 =-\begin{bmatrix}A(\lambda)v\\0\end{bmatrix},\qquad \|v\|_2=1.
\end{equation}
The next iterate uses $\lambda+\delta\lambda$ with damping inside the cell.
Without a supplied derivative, four-direction finite differences are used
only for refinement; the trace cross-check is explicitly marked unavailable.
Singular vectors supply left/right residual checks, not the analytic target
of a scalar zero search. Vectors are mapped back as $v=R_0\widetilde v$ and
$w=L_0^*\widetilde w$ and normalized in the original coordinates.

\section{Multiplicity, clusters and numerical limits}

For $K=(s+1)I_2$, the algebraic count and nullity are both two; the projected
derivative is nonsingular, supporting a numerical semisimple double mode.
For
\begin{equation}
 K=\begin{bmatrix}s+1&1\\0&s+1\end{bmatrix},
\end{equation}
the count is two but nullity is one. The initial engine returns a counted
cluster with unknown local multiplicity, not an invented simple mode or a
computed Jordan chain. Very close distinct roots may also remain clusters.
\code{countComplete} can therefore be true while \code{complete} is false.
The returned representative must be read together with \code{localCounts},
\code{multiplicity}, \code{nullity} and \code{clusters}.

Default local confirmation squares have half-width
$10^{-8}(1+|\lambda|)$; rank decisions use $10^{-8}$ times the fixed scaled
reference norm. Residuals use that same regular-reference norm rather than
the vanishing matrix norm at a mode. Unscaled residuals are also returned.
No partial multiplicities, contour moments or sparse large-scale extraction
are claimed. Subdivision is not limited to the matrix dimension: the scalar
matrix $\sinh s$ already has seven roots in $[-1,1]\times[-10,10]$.

\section{Determinant and continuous-wing comparisons}

For $K=c\left[\begin{smallmatrix}0&s+2\\s+1&0\end{smallmatrix}\right]$,
$c=10^{-200}$ or $10^{200}$, the raw determinant at $s=0.5$ is zero or
negative infinity in floating-point arithmetic. The scalar determinant
search is unresolved; the matrix search finds $-1,-2$ with complete counts.
This demonstrates a representable-range advantage, not immunity to
ill-conditioning or to overflow inside the supplied function.

For the wing, search the implicit boundary-value matrix of
Chapter~\ref{ch:wing}, not the meromorphic hybrid port matrix of
Chapter~\ref{ch:hybrid}. Keep the number of exact numerical pieces fixed
during each search so that the dimension does not change:
\begin{lstlisting}
wing = wing_model('goland'); U = 150; pieces = 2;
K = @(s) wing_matrix(s,U,wing,pieces);
box = [-40 15 1 320]; % above the Theodorsen branch cut
[r,i] = cmodes(wing_cdyn(U,wing,pieces),box, ...
    'AssumeAnalytic',true);
[d,j] = croots(@(s) det(full(K(s))),box, ...
    'AssumeAnalytic',true);
\end{lstlisting}

\begin{table}[htbp]\centering\small
\caption{Matrix search versus determinant roots and independent dry
frequencies. Error is $|\lambda-\lambda_{\rm ref}|/(1+|\lambda_{\rm ref}|)$.}
\begin{tabular}{llrrrr}\toprule
Model & Pieces & Dimension & Count & Determinant error & Dry exact error\\\midrule
Dry & 1 & 12 & 4 & $2.25\,10^{-13}$ & $2.25\,10^{-13}$\\
Dry & 2 & 18 & 4 & $2.82\,10^{-15}$ & $5.63\,10^{-16}$\\
Dry & 4 & 30 & 4 & $2.60\,10^{-15}$ & $4.22\,10^{-16}$\\
150 m/s & 1 & 12 & 3 & $5.88\,10^{-16}$ & ---\\
150 m/s & 2 & 18 & 3 & $9.21\,10^{-16}$ & ---\\
150 m/s & 4 & 30 & 3 & $2.50\,10^{-16}$ & ---\\\bottomrule
\end{tabular}
\end{table}

\begin{figure}[htbp]\centering
\includegraphics[width=\textwidth]{matrix_modes_comparison.png}
\caption{Continuous-wing modes from the matrix and determinant routes;
agreement across fixed exact subdivisions.}
\end{figure}

All six runs have reconciled counts. The dry search window is
$[-2,2]\times[1,320]$; the air-loaded window is the one in the listing
(its left edge is kept well away from the mode near $-29.6+57i$).
Subdivision changes the numerical representation, not the constant-strip
equations. The improved dry frequency error with two/four pieces supports
a conditioning benefit in this example, not a universal monotonic rule.
Matrix/determinant agreement is not independent validation of the shared
aerodynamic model; the physical checks remain those of Chapter~\ref{ch:wing}.
The reproducible script is
\path{examples/continuous_wing/run_matrix_modes_comparison.m}.

\section{The flutter mode shape}

The right null vector at a mode holds the scaled states at every node of
$K$. Propagating the root state $z_0$ with the exact strip propagators,
$z(y)=e^{A(\lambda,U)y}z_0$ within each strip, gives the mode shape at
every span position without interpolation or modal basis
(\code{wing\_mode\_shape}). At the flutter speed $U_f=136.984$~m/s,
\code{cmodes} returns the mode $\lambda=70.0330i$, and the propagated shape
satisfies the free-tip conditions to $5\times10^{-14}$ relative to the root
loads (Figure~\ref{fig:fluttermode}). At the tip the twist is as large as
the bending motion, $b|\alpha|/|w|=1.034$, and lags it by $62^\circ$. This
phase difference lets the air do net work on the structure in each cycle,
the mechanism of bending--torsion flutter. One and two exact pieces give
the same ratio and phase.

\begin{figure}[htbp]\centering
\includegraphics[width=\textwidth]{wing_flutter_mode.png}
\caption{Flutter mode of the continuous Goland wing. Left: bending and
twist amplitudes along the span (tip deflection normalized to one).
Right: tip motion over one cycle.}
\label{fig:fluttermode}
\end{figure}

This qualifies the contour-based M3 migration, not general MIMO spectra.
Adequate sampling and analyticity remain assumptions; \code{certified} is
always false. Work limits do not control allocations inside callbacks.
Full contracts and deferred capabilities are documented in
\path{docs/api/cmodes.md}.
